% GENERATED FILE. Edit canonical lessons, then run npm run generate:books.
% canonical-source-hash: fnv1a64:4a45adedc5c9df8e
% canonical-lessons: PA-W16-form-mixed-return, PA-W16-form-timed-return, PA-W16-message-pacing-return

\chapter{Return to Three Familiar A1 Writing Moves}
\label{ch:pa-three-familiar-a1-writing-moves}
\hlchaptermodality{162}

\begin{chapteropening}
\textbf{By the end of this chapter:} \emph{I can retrieve mixed-card form choices, preserve a six-field timed first attempt, and pace a named-reader message opening without treating these short returns as complete A1 mocks.}

Inspect a preserved mixed-card timed form attempt, then write a fresh, short named-reader message opening without model or mock credit.
\end{chapteropening}

\section[\texorpdfstring{\pa{ਨਾਂ}: · \pa{ਭਾਸ਼ਾ} … (return to three …)}{ਨਾਂ: · ਭਾਸ਼ਾ … (return to three …)}]{Return to three old card choices}

\label{lesson:PA-W16-form-mixed-return}



\begin{quote}
The old A · \pa{ਕ} and B · \pa{ਖ} cards choose values you have already learned. For this supported return, reopen the old word banks: A has \textbf{\pa{ਅਮਨ}}, \textbf{\pa{ਪੰਜਾਬੀ}}, \textbf{\pa{ਪਿੰਡ}}; B has \textbf{\pa{ਮਨਨ}}, \textbf{\pa{ਹਿੰਦੀ}}, \textbf{\pa{ਸ਼ਹਿਰ}}. No new word or person is being introduced.
\end{quote}

\subsection*{Writing — supported mixed selection}
With the banks visible, write the old value chosen by each mark. Go one line at a time. There is no clock and no independent-writing credit here.

\begin{quote}
\pa{ਨਾਂ} (B · \pa{ਖ}): \_\_\_\_\_\_
\end{quote}

\begin{quote}
\pa{ਭਾਸ਼ਾ} (A · \pa{ਕ}): \_\_\_\_\_\_
\end{quote}

\begin{quote}
\pa{ਰਿਹਾਇਸ਼} (B · \pa{ਖ}): \_\_\_\_\_\_
\end{quote}

\subsection*{Your turn — check after writing}
First check the three card choices: \textbf{\pa{ਮਨਨ}}, \textbf{\pa{ਪੰਜਾਬੀ}}, \textbf{\pa{ਸ਼ਹਿਰ}}. On a separate pass, check spelling and the space after each colon. Close the banks again. This small supported return is not a complete A1 mock.

\begin{tcolorbox}[breakable,colback=teal!4,colframe=teal!35!black,title={Before you move on}]
Say which card mark belonged to each of your three lines. The next return uses all six familiar fields without a visible bank.
\end{tcolorbox}
\section[\texorpdfstring{\pa{ਨਾਂ}: · \pa{ਭਾਸ਼ਾ} … (return to a …)}{ਨਾਂ: · ਭਾਸ਼ਾ … (return to a …)}]{Return to the six familiar fields}

\label{lesson:PA-W16-form-timed-return}



\begin{quote}
Close the value banks and earlier forms. Have a blank page and a three-minute timer ready. The marks below select only old values. Do not look at a model while the timer runs.
\end{quote}

\subsection*{Writing — independent timed form}
Start the timer. Write all six marked fields in Gurmukhi; use Gurmukhi digits for age and phone. Stop at three minutes, leaving unfinished lines as they are. Keep this first attempt unchanged.

\begin{quote}
\pa{ਨਾਂ} (B · \pa{ਖ}): \_\_\_\_\_\_
\end{quote}

\begin{quote}
\pa{ਭਾਸ਼ਾ} (A · \pa{ਕ}): \_\_\_\_\_\_
\end{quote}

\begin{quote}
\pa{ਰਿਹਾਇਸ਼} (B · \pa{ਖ}): \_\_\_\_\_\_
\end{quote}

\begin{quote}
\pa{ਕੰਮ} (A · \pa{ਕ}): \_\_\_\_\_\_
\end{quote}

\begin{quote}
\pa{ਉਮਰ} (B · \pa{ਖ}): \_\_\_\_\_\_
\end{quote}

\begin{quote}
\pa{ਫ਼ੋਨ} (A · \pa{ਕ}): \_\_\_\_\_\_
\end{quote}

\subsection*{Your turn — repair separately}
After stopping, reopen the old banks and check values, spelling, spacing, digit order, and phone grouping one at a time. Repair on another page; it does not replace the clocked evidence. This return is not a complete A1 mock.

\begin{tcolorbox}[breakable,colback=teal!4,colframe=teal!35!black,title={Before you move on}]
Record how many lines were complete at the stop. The full form-plus-message condition belongs to the later two complete mocks.
\end{tcolorbox}
\section[(named reader, short clock)]{Return to the reader and the clock}

\label{lesson:PA-W16-message-pacing-return}



\begin{quote}
Keep the previous lesson's six-field first attempt. In one minute, point to a line whose card mark you selected correctly and a line you would repair. Do not rewrite the timed page; its first-attempt evidence stays as it was.
\end{quote}

\subsection*{Writing — one small timed opening}
Close the earlier message pages. As Aman, write to Manan in Gurmukhi. In ninety seconds, greet your named reader and ask one familiar wellbeing question. Choose your own taught words; no sample sentence is visible. Stop on time and keep the first attempt.

\subsection*{Your turn — inspect one boundary}
After the timer, check whether Manan is clearly the reader. Then check the question mark and spacing. A repair belongs on another page and is not timed evidence. This opening is not a complete A1 mock.

\begin{tcolorbox}[breakable,colback=teal!4,colframe=teal!35!black,title={Before you move on}]
The complete paper still needs a fresh purpose, a 30–40-word message, and a twenty-minute form-plus-message clock. Those conditions belong to the two separate mocks, not this short return.
\end{tcolorbox}
